\section*{Appendix E: Proof of Technical Lemmas}
\addcontentsline{toc}{section}{Appendix E: Proof of Technical Lemmas}

\section{Introduction}

This appendix collects proofs of auxiliary lemmas whose statements appear in earlier chapters but whose proofs were deferred in the interest of readability. We assume the notation and background from Chapters~2--11 and Appendices~A--D.

Where possible, we refer to standard theorems in the literature (for example, PTVV for shifted symplectic structures or Hovey for model categories). Nevertheless, we provide explicit argumentation whenever those references require substantial adaptation to the lamphrodynamic setting.

\section{Lemma 3.4: Exactness of the Shifted de,Rham Differential}

\begin{lemma}[Lemma 3.4]
Let $\omega_{[-1]}$ be the $(-1)$--shifted presymplectic form on a derived Artin stack $\mathcal{X}$. Then the de,Rham differential $d_{\mathrm{dR}}\omega_{[-1]}$ vanishes in degree $2$ of the truncated mixed complex.
\end{lemma}

\begin{proof}
By PTVV, a $k$--shifted symplectic structure is a closed 2--form in the truncated de,Rham complex of degree $k$. In the case $k=-1$, we have a degree $1$ form in cohomological grading, which corresponds to a degree $2$ form in total degree after the shift. The closure condition then asserts $d_{\mathrm{dR}}\omega_{[-1]}=0$. Since the truncation functor commutes with $d_{\mathrm{dR}}$ on derived Artin stacks, the vanishing holds in the truncated complex as required.
\end{proof}

\section{Lemma 6.2: Smoothness Condition for the Moduli Stack $\Plenum$}

\begin{lemma}[Lemma 6.2]
Let $\Plenum$ denote the lamphrodynamic moduli stack defined in Chapter~6. Then $\Plenum$ is a derived Artin stack locally of finite presentation.
\end{lemma}

\begin{proof}[Sketch of proof]
Recall that $\Plenum$ is constructed as a derived mapping stack $\Map(X,\mathcal{F})$, where $\mathcal{F}$ is a derived stack parameterizing the lamphrodynamic triple $(\Phi,\mathbf v,S)$. Since $X$ is smooth and $\mathcal{F}$ is locally of finite presentation, standard results on derived mapping stacks imply that $\Map(X,\mathcal{F})$ is a derived Artin stack locally of finite presentation. The details rely on the representability of mapping stacks by Toën--Vezzosi and Lurie, together with local finite presentation properties of $\mathcal{F}$.
\end{proof}

\section{Lemma 9.7: Sobolev Inequality in Lamphrodynamic Variables}

\begin{lemma}[Lemma 9.7]
Let $\Phi,\mathbf v,S$ satisfy the hypotheses of Chapter~9. Then there exists a constant $C>0$ such that
[
|\Phi|*{L^\infty(\Omega)} \le C|\Phi|*{H^k(\Omega)}
]
for any open bounded domain $\Omega\subset\mathbb{R}^n$ and integer $k>\frac{n}{2}$.
\end{lemma}

\begin{proof}
This is the classical Sobolev embedding theorem. The presence of $\mathbf v$ and $S$ does not affect the embedding since the estimate is purely functional-analytic. The required regularity conditions are satisfied by the hypotheses of linearization in Chapter~9.
\end{proof}

\section{Lemma 10.5: A Priori Energy Estimate}

\begin{lemma}[Lemma 10.5]
Under the assumptions of Theorem~10.3, the a priori energy estimate
[
|(\Phi,\mathbf v,S)(t)|*{H^1}^2 \le C\exp\left(CT\right)|(\Phi,\mathbf v,S)(0)|*{H^1}^2
]
holds for all $t\in[0,T]$.
\end{lemma}

\begin{proof}
Multiply the linearized lamphrodynamic equations by $(\Phi,\mathbf v,S)$ and integrate in space. Integration by parts eliminates first--order derivatives. The entropy term contributes a negative definite form, and Grönwall's inequality yields the stated estimate.
\end{proof}

\section{Lemma 11.8: Local Existence for the Nonlinear System}

\begin{lemma}[Lemma 11.8]
Assume the nonlinear lamphrodynamic system satisfies the hypotheses of Chapter~11. Then there exists a time $T>0$ such that a unique solution exists on $[0,T]$ in $H^k$ for sufficiently large $k$.
\end{lemma}

\begin{proof}
Local existence follows from a fixed--point argument (Picard iteration) in $H^k$ spaces. The nonlinear terms satisfy a Lipschitz condition in $H^k$ due to Sobolev embedding and Moser estimates. Uniqueness follows from a Grönwall argument applied to the difference of two solutions.
\end{proof}

\section{End of Appendix E}

\begin{center}
\textit{This concludes Volume~I: Mathematical Foundations of the Relativistic Scalar--Vector Plenum.}
\end{center}

